\chapter{Анализ предметной области}


\section{Определение рекомендательной системы}

Рекомендательная система — это программа, предоставляющая рекомендации на основе данных о пользователях
и о приобретаемых ими предметах~\cite{def_rec}. Другими словами, такие программы могут моделировать возможный будущий
выбор покупателей. Такие системы, обрабатывая полученную информацию, предполагают, какой именно
объект вызовет наибольший интерес у пользователя. Рекомендательная система включает в себя цикл действий по работе с
данными, начиная от получения сырых данных о пользователе, объектах и предпочтениях пользователя и заканчивая
предоставлением рекомендации потребителю.


%\section{Типы рекомендательных систем}
%
%\subsection{Коллаборативная фильтрация (Collaborative Filtering)}
%
%В этом типе рекомендательных систем фильтрация элементов из большого набора
%альтернатив осуществляется совместными усилиями схожих пользовательских
%предпочтений. Основная концепция заключается в том, что если два
%пользователя имели одинаковые интересы в прошлом, они могут иметь такие
%же интересы и в будущем.~\cite{obz_sovr} Таким образом, система рекомендует пользователю те элементы,
%которые нравятся «похожим» пользователям.
%
%\subsection{Контентно-ориентированные (Content-based)}
%
%В данном типе рекомендательных систем рекомендации формируются на основе сходства свойств объектов, которые пользователь уже оценил,
%с другими объектами. Система анализирует характеристики контента (жанры фильмов, тематику статей,
%технические параметры товаров) и подбирает похожие объекты.
%
%\subsection{Основанные на знаниях (Knowledge-based) рекомендации}
%
%В данном типе система полагается на явные знания о домене, предпочтениях пользователя и особенностях объектов.
%Рекомендации строятся на основании правил и логических заключений, онтологий, семантических сетей или экспертных систем.
%Такие системы могут включать взаимодействие с пользователем для определения его потребностей, задавая вопросы о
%значимости тех или иных характеристик. Однако важно, чтобы это взаимодействие было удобным и минимизировало
%когнитивную нагрузку, особенно при большом числе параметров.~\cite{rek_intro}
%
%\subsection{Гибридные рекомендательные системы (Hybrid Systems)}
%
%Гибридные системы объединяют два и более методов (коллаборативный, контентный, основанный на знаниях)
%для компенсации слабых сторон и усиления сильных.~\cite{rek_intro}
%
%\clearpage

\section{Основные вехи развития}

История развития рекомендательных систем начинается с конца
1990-х годов, с появлением цифрового видеомагнитофона TiVo. В нем
использовалась модель коллаборативной фильтрации.

В конце 1990-х -- начале 2000-х Amazon и другие о Крупные коммерческие игроки начали применять
рекомендации для увеличения продаж, внедряя методы фильтрации, чтобы в реальном времени предлагать покупателям
товары, похожие на уже просматриваемые или купленные.

К середине 2000-х появились фундаментальные обзоры, такие как работа <<Adomavicius \& Tuzhilin>>~\cite{hist} (2005),
систематизировавшая знания в области рекомендаций. Формируются чёткие таксономии методов,
а также определяются стандартные метрики для оценки качества.

Компания Netflix в 2006 году организовала соревнование
Netflix Prize. В течение трех лет команды разработчиков добились
только прироста точности в, но эти решения не были применимы
на практике, так как требовали огромных вычислительных ресурсов.
В результате соревнования NetflixPrize выиграл гибридный
подход для систем рекомендаций.

В 2010-е годы, с развитием машинного обучения возникает интерес к контекстуальным
рекомендациям (учёт времени, места, устройства) и к сложным доменам, таким как рекомендации
туристических направлений или образовательного контента.

